\chapter{Operations, Team and Roadmap}\label{ch:ops}

\section{Operating model}\label{sec:opmodel}
ClosetTwin runs as three parts: a capture widget on the retailer's product page, an inference backend that turns the scan, the garment data and the stated fit preference into a size, and a data feed that returns size and return signals to the retailer. The backend runs on rented cloud capacity, so cost grows with scans and not with fixed equipment. We assume a cost to serve of \usd{0.0075} per order, which is a scan cost of \usd{0.03} on the 25\% of orders that include a scan (both \assumed{assumed}).

\section{Team and roles}\label{sec:team}
The team grows from 3 people in year 1 to 4, 5, 9 and 12 in years 2 to 5 (\assumed{assumed}, at a fully loaded cost of \usd{90,000} per person-year). The first three roles are engineering, data science and partnerships with a compliance duty attached. The ninth to twelfth positions arrive in year 4 with the Europe entry and add sales, integration support and privacy compliance. The cost of \usd{90,000} per person-year is a blended assumption with no benchmark behind it, for Israeli engineers and European salespeople alike. The model's operating cost is this payroll plus three assumed lines for marketing and demos, cloud and tooling, and legal and compliance (Chapter~\ref{ch:finance}); travel and general administration are not modelled (Section~\ref{sec:limits}). We describe roles only and name no individuals. The size is deliberately lean: Israeli revenue cannot carry a larger team, as Chapter~\ref{ch:finance} shows.

\section{Roadmap}\label{sec:roadmap}
Figure~\ref{fig:f4} shows six workstreams with relative quarters, all \assumed{assumed}. The minimum viable product is built in year 1, quarters 1 and 2. The compliance milestone for Israel's Privacy Protection Law as amended by Amendment 13~\cite{il_amend13_2024} (privacy review, consent flow and a check of the database-notification duty) falls in year 1, quarter 2, before any pilot data are collected. The first Israeli pilots run in year 1, quarters 3 and 4, from warm design-partner contacts, because a nine-month sales cycle leaves no room for cold starts. The first paying Israeli clients arrive in year 2. A review under Article 9 of the GDPR~\cite{gdpr2016} follows in year 3, quarter 3. Europe entry, meaning pilots, starts in year 4, and the first paying European clients arrive in year 5.

\begin{figure}[tbp]
% F4: roadmap Gantt strip, 6 rows, relative quarters, all assumed (A18, A25, A26, A45).
% No caption/label here: the chapter owns them (fig:f4). Quarter q spans [q,q+1] at 0.55 cm.
\begin{tikzpicture}[x=0.55cm,y=0.8cm]
  \foreach \y in {1,...,5}{
    \pgfmathtruncatemacro{\q}{(\y-1)*4}
    \draw[ct guide] (\q,0.4) -- (\q,-6);
    \node[ct label,anchor=south west] at (\q,0.4) {Y\y};
  }
  \draw[ct guide] (20,0.4) -- (20,-6);
  \node[ct label,anchor=east] at (-0.3,-0.5) {MVP};
  \node[ct label,anchor=east] at (-0.3,-1.5) {Amendment 13 compliance};
  \node[ct label,anchor=east] at (-0.3,-2.5) {Israeli pilots};
  \node[ct label,anchor=east] at (-0.3,-3.5) {Israeli paying clients};
  \node[ct label,anchor=east] at (-0.3,-4.5) {GDPR Art.\ 9 review};
  \node[ct label,anchor=east] at (-0.3,-5.5) {Europe: pilots, paying};
  \node[ct bar,anchor=west,minimum width=1.1cm] at (0,-0.5) {};
  \node[ct bar milestone] at (1.5,-1.5) {};
  \node[ct bar,anchor=west,fill=pink,minimum width=2.2cm] at (2,-2.5) {};
  \node[ct bar,anchor=west,fill=mochadeep,minimum width=8.8cm] at (4,-3.5) {};
  \node[ct bar milestone] at (10.5,-4.5) {};
  \node[ct bar,anchor=west,fill=pink,minimum width=2.2cm] at (12,-5.5) {};
  \node[ct bar,anchor=west,fill=mochadeep,minimum width=2.2cm] at (16,-5.5) {};
\end{tikzpicture}

\caption{Roadmap, years 1 to 5 (relative quarters, all assumed)}\label{fig:f4}
\end{figure}

\section{Dependencies}\label{sec:dependencies}
Three dependencies matter. The first is the pilot result (Section~\ref{sec:pilot-design}). The second is a consent flow that holds up legally under Amendment 13. The third is the 50\% conversion of pilots to paying clients, which stays untested until the first pilots finish. Europe's revenue depends on all three, and Chapter~\ref{ch:risks} rates the matching risks.
\ctchapterend
